\documentclass[11pt]{article}
\usepackage[a4paper,margin=28mm]{geometry}
\usepackage{amsmath,amssymb,amsthm,mathtools,hyperref}
\newcommand{\ddc}{dd^c}
\newtheorem{theorem}{Theorem}
\newtheorem{lemma}[theorem]{Lemma}
\theoremstyle{remark}\newtheorem{remark}[theorem]{Remark}
\title{Endpoint cost for all minimizing Gauduchon metrics\newline on nonclosed warped paths}
\author{}\date{}
\begin{document}\maketitle
\noindent\textbf{Working draft.} This new root-developed candidate has not
passed complete native review, independent reconstruction or final comparison.
\begin{abstract}
We propose an exact endpoint-cost theorem on a compact K\"ahler manifold
times a curve. A smooth nonnegative vertical coefficient determines
nonclosed paths with actual balanced minimizers. The cost is independent
of every choice of attaining Gauduchon metric, and its improper integral
equals the drop in root volume. Tubular codimension and vanishing order
give five quantitative integrable profiles. Positive endpoint volume is
compatible with failure of bigness, which is excluded for every genuine
quasi-psh strict representative by a smooth closed test.
\end{abstract}

Use $\ddc=\frac{i}{2\pi}\partial\bar\partial$ and ordinary wedge powers.
Let $(B,\alpha)$ be a connected compact K\"ahler $m$-fold, $m\geq2$,
$W=\int_B\alpha^m>0$. Let $C$ be a connected compact complex curve
with a smooth positive probability form $w$, and $\sigma$ a smooth
positive probability on $B$. Put $X=B\times C$, $n=m+1$,
$\mu=\sigma\wedge w$; suppress pullbacks. Fix $a\geq0$ smooth and $t>0$.
Consider the strict path and nef endpoint
\[
 \theta_s=\alpha+(a+t-s)w,\quad 0\leq s<t,
 \qquad\theta_t=\alpha+aw,\qquad\eta=w.
\]
For nonconstant $a$, the torsion is $d\theta_s=da\wedge w\ne0$.
Dependence on the specified $\sigma$ remains in the formulas below.

For every smooth positive Gauduchon metric define
\[
 D_\mu(\gamma)=\int_X(\mu/\gamma^n)^{1/n}\gamma^n,
 \quad\lambda_\mu(\theta)=
 \inf_{\ddc\gamma^m=0}\frac{\int_X\theta\wedge\gamma^m}{D_\mu(\gamma)},
 \quad V_\mu=\lambda_\mu^n.
\]
The endpoint is defined by decreasing Hermitian regularization, agreeing
with the nonnegative full infimum as proved below. A background is big
if it has a genuine global quasi-psh potential making it dominate a
positive multiple of a Hermitian metric.

\section{Actual minimizers and the full volume}
Write $\epsilon=t-s$, $L=a+\epsilon$, and
\[
 H(\epsilon)=\int_B\frac{\sigma}{a+\epsilon},\qquad
 J(\epsilon)=\int_B\frac{\sigma}{(a+\epsilon)^2}.
\]
Smooth Monge--Amp\`ere existence in the closed positive class
$[\alpha]$, in the formulation of \cite[Corollary 1]{TW}, gives
\[
 A_\epsilon=\alpha+\ddc\phi_\epsilon>0,\qquad
 A_\epsilon^m=\frac{W}{H(\epsilon)}\frac{\sigma}{L}.
\]
The target has total mass $W$, so integration in the closed class fixes
the source scalar. Our potential is $2\pi$ times the source's
$i\partial\bar\partial$ potential. No endpoint equation or uniform
endpoint estimate is imported.

Set $\Omega_\epsilon=A_\epsilon+Lw$,
$c_\epsilon=nW/H(\epsilon)$, and
$\gamma_s=L^{-1/m}\Omega_\epsilon$. Then
\[
 \Omega_\epsilon^n=c_\epsilon\mu,
 \qquad \gamma_s^m=L^{-1}A_\epsilon^m+mA_\epsilon^{m-1}\wedge w.
\]
Its first cofactor summand is a base top-degree form and its second is
closed, so $d\gamma_s^m=0$. For every admissible test $\xi$, smooth
Stokes and matrix AM--GM give
\[
 \int_X\theta_s\wedge\xi^m=\int_X\Omega_\epsilon\wedge\xi^m
 \geq c_\epsilon^{1/n}D_\mu(\xi).
\]
For the named $\gamma_s$ the numerator is $nW$ and denominator
$c_\epsilon^{m/n}H(\epsilon)$; hence it attains equality. In particular
\begin{equation}\label{eq:volume}
 V_\mu(\theta_s)=\frac{nW}{H(t-s)}.
\end{equation}
This is the full infimum over all smooth positive Gauduchon metrics,
not an infimum over product tests.

\section{Cost independent of the minimizing selection}
For an actual attaining metric $\xi$ at $s$ define
\[
 C_\xi(s)=\frac{\int_X w\wedge\xi^m}{D_\mu(\xi)}.
\]
The named balanced metric satisfies
\[
 \int_X w\wedge\gamma_s^m=\frac{WJ}{H},
 \quad D_\mu(\gamma_s)=(nW/H)^{m/n}H,
 \quad C(s)=n^{-m/n}W^{1/n}H^{-(n+1)/n}J.
\]
The last expression is also $-\lambda_\mu'(\theta_s)$ by
differentiating \eqref{eq:volume}.

\begin{lemma}\label{lem:all}
Every actual minimizing Gauduchon metric at $s$ has cost $C(s)$.
No differentiability of a selected family of metrics is required.
\end{lemma}
\begin{proof}
Hold any such $\xi$ fixed. The all-test infimum at a nearby path point
gives, for both signs of sufficiently small $\delta$,
\[
 \lambda_\mu(\theta_{s+\delta})
 \leq\lambda_\mu(\theta_s)-\delta C_\xi(s).
\]
The explicit volume is differentiable on the strict path. Dividing by
positive and negative $\delta$ and taking their limits forces
$C_\xi(s)=-\lambda_\mu'(\theta_s)=C(s)$.
At $s=0$ the same argument uses the short strict extension to negative
$s$, for which $a+t-s$ is still positive.
\end{proof}

\begin{theorem}\label{th:endpoint}
For any smooth $a\geq0$ and any selection of actual attaining metrics,
the improper integral exists and
\[
 \int_0^t C(s)\,ds
 =\lambda_\mu(\theta_0)-\lambda_\mu(\theta_t)<\infty,
 \qquad V_\mu(\theta_t)=\frac{nW}{H(0)}.
\]
Here $H(0)=\int_B\sigma/a$, with endpoint volume zero if it is infinite.
\end{theorem}
\begin{proof}
The full root infimum at the nef endpoint is the limit of its strict
path values. Monotonicity gives one inequality. For any fixed test its
quotient converges to its endpoint quotient; taking the infimum gives
the reverse inequality. The same argument with any positive Hermitian
regularization identifies the original endpoint definition.
Monotone convergence yields $H(\epsilon)\to H(0)$, so
\eqref{eq:volume} gives the asserted actual endpoint value.

On each compact strict path segment the fundamental theorem gives
\[
 \int_0^r C(s)\,ds=\lambda_\mu(\theta_0)-\lambda_\mu(\theta_r).
\]
The integrand is nonnegative and the endpoint limit just proved is
finite. Letting $r\uparrow t$ proves the improper identity.
Neither integrability nor an endpoint Monge--Amp\`ere solution has been
assumed. Lemma~\ref{lem:all} covers every minimizing selection.
\end{proof}

\section{Quantitative geometric profiles}
Suppose $\{a=0\}$ is a smooth compact real submanifold of codimension
$d>0$ and in a fixed tube
$a\asymp\operatorname{dist}(\cdot,\{a=0\})^k$, $k>0$,
with positive two-sided constants. Assume $a$ is bounded below outside
the tube, and put $r=d/k$. These are geometric hypotheses on the
coefficient, not hypotheses on the desired cost.

Tubular integration compares $H$ and $J$ to
\[
 \int_0^R\frac{u^{d-1}}{u^k+\epsilon}\,du,
 \qquad \int_0^R\frac{u^{d-1}}{(u^k+\epsilon)^2}\,du.
\]
Splitting at $\epsilon^{1/k}$ gives the thresholds one and two. Therefore
\begin{equation}\label{eq:profiles}
 C(t-\epsilon)\asymp
 \begin{cases}
 \epsilon^{(1-r)/n-1},&r<1,\\
 \epsilon^{-1}\bigl[\log(1/\epsilon)\bigr]^{-(n+1)/n},&r=1,\\
 \epsilon^{r-2},&1<r<2,\\
 \log(1/\epsilon),&r=2,\\
 1,&r>2.
 \end{cases}
\end{equation}
In the final case the cost tends to the finite positive value
$n^{-m/n}W^{1/n}H(0)^{-(n+1)/n}J(0)$ by dominated convergence.
Each profile is integrable: the first power is greater than $-1$;
the logarithmic power at $r=1$ exceeds one; the power for $1<r<2$
is again greater than $-1$. The endpoint volume is positive exactly
when $d>k$. With finitely many disjoint finite-order components the
same comparisons use their smallest codimension/order ratio.

\section{Bigness and an explicit positive-volume nonbig endpoint}
If $a$ has a zero, $\theta_t$ is not big. Indeed suppose for a genuine
quasi-psh $u$ that
\[
 \theta_t+\ddc u\geq\delta(\alpha+w),\qquad\delta>0.
\]
Test against the pullback of an arbitrary smooth nonnegative base
top-degree form $\zeta$. This test is $\ddc$-closed by dimension, so
distributional Stokes gives
$\int_B(a-\delta)\zeta\geq0$.
A nonzero $\zeta$ supported where $a<\delta$ contradicts this.
The argument uses local integrability of genuine quasi-psh functions
and smooth tests; no singular potential is restricted to a curve fiber.
If $a>0$ everywhere then $\theta_t$ itself is positive and big.
Thus in this family bigness is exactly the absence of zeros.

For a concrete genuinely nonclosed example take $B=\mathbb P^m$,
$\alpha=\kappa$ normalized Fubini--Study, $\sigma=\kappa^m$,
$C=\mathbb P^1$, and
\[
 a([Z])=\frac{\sum_{j=1}^m|Z_j|^2}{\sum_{j=0}^m|Z_j|^2}.
\]
Its single zero is quadratic in a real codimension-$2m$ affine chart,
so $r=m>1$. The law of $a$ has density $m x^{m-1}dx$ on $[0,1]$:
polar integration in the $Z_0$ chart gives radial density
$m q^{m-1}(1+q)^{-m-1}dq$, and $x=q/(1+q)$ gives this law.
Hence $W=1$, $H(0)=m/(m-1)$ and
\[
 V_\mu(\theta_t)=\frac{(m+1)(m-1)}m>0,
\]
although the endpoint is not big. Its cost is logarithmic at $m=2$,
and has a finite positive limit for $m>2$.
These statements retain the nonclosed torsion and the chosen probability.

\section{Closed Morse comparison and the closed-core boundary}
In the closed subcase $a\equiv a_0\geq0$, put
$P=\alpha+(a_0+t)w$ and $Q=tw$ on the compact K\"ahler product.
Both forms are $d$-closed and their real $(1,1)$ classes are nef.
Since $w^2=0$ and $\alpha^{m+1}=0$,
\[
 \int_X P^n=nW(a_0+t),\qquad
 \int_X P^{n-1}\wedge Q=tW.
\]
The cohomological nef Morse inequality
\cite[Theorem 1.1]{Morse} therefore reads
\[
 \operatorname{vol}([P]-[Q])\geq
 \int_XP^n-n\int_XP^{n-1}\wedge Q
 =nWa_0=V_\mu(\theta_t).
\]
For $a_0>0$, the final equality also follows directly from
$H(0)=1/a_0$; for $a_0=0$ both displayed values are zero.
The expression is independent of $\sigma$ because $\int_B\sigma=1$.
Here $\operatorname{vol}$ is the usual volume of a closed cohomology
class, extended by zero outside the pseudoeffective cone.
This comparison uses the established closed-class theorem, without
proposing another proof of its general assertion.

For nonconstant $a$, $d\theta_t=da\wedge w\ne0$ and
$H(0)=\int_B\sigma/a$ retains the specified density.
There is no class $[\theta_t]$ to which that closed-class theorem applies.
In particular the preceding comparison supplies no nonclosed
intersection-number formula.

Moreover every endpoint with a zero lies outside the closed-core
hypothesis. Suppose, even up to a smooth global potential gauge, that
$\theta_t=\chi+\beta+\ddc f$, where $\chi$ is smooth real and
$d$-closed, $[\chi]$ is nef with $\int_X\chi^n>0$, and
$\beta\geq0$ is smooth. By \cite[Theorem 2.12]{DP}, the class
$[\chi]$ contains a K\"ahler current. The K\"ahler
$\partial\bar\partial$ lemma and the rescaling of potentials give a
genuine global quasi-psh $v$ with
$\chi+\ddc v\geq\delta(\alpha+w)$ for some $\delta>0$.
Then
\[
 \theta_t+\ddc(v-f)=\chi+\ddc v+\beta
 \geq\delta(\alpha+w),
\]
contradicting the nonbigness proof for every genuine potential.
This excludes every such closed nef positive-volume core, rather than
only a chosen product core.

The full product volume calculation, tubular inverse-moment estimates
and a scaled version of the projective positive-volume nonbig example
already occur in the companion manuscript and the checked warped-cone
calculation \cite{Companion}. They serve here as context.
The present theorem gives the exact improper endpoint identity for
every attaining selection on an arbitrary compact K\"ahler base,
and the five geometric cost profiles. No literature priority or
classification of arbitrary nonclosed paths is asserted.

\begin{remark}
This theorem supplies a geometric endpoint family beyond conclusions
requiring a closed positive-volume core. It does not classify arbitrary
paths or assert that positive nonclosed root volume implies bigness.
Only the closed base mass $W$ is invariant under smooth potential
changes; no nonclosed intersection-number invariance, metric uniqueness,
endpoint PDE solvability or literature priority is asserted.
\end{remark}
\begin{thebibliography}{9}
\bibitem{TW} V. Tosatti and B. Weinkove, \emph{The complex Monge--Amp\`ere
equation on compact Hermitian manifolds}, arXiv:0910.1390v2
(11 November 2009), Corollary 1 and closed-background normalization,
pp.~1--2; \url{https://arxiv.org/abs/0910.1390}.
\bibitem{DP} J.-P. Demailly and M. P\u{a}un,
\emph{Numerical characterization of the K\"ahler cone of a compact
K\"ahler manifold}, arXiv:math/0105176v2 (24 September 2001),
Theorem 2.12; \url{https://arxiv.org/abs/math/0105176}.
\bibitem{Morse} V. Tosatti, \emph{Transcendental Morse inequality on
K\"ahler manifolds}, arXiv:2608.30153v2 (8 September 2026),
Theorem 1.1; \url{https://arxiv.org/abs/2608.30153}.
\bibitem{Companion} M. Xia, \emph{Non-K\"ahler pluripotential theory},
companion repository: the radial positive-volume nonbig example,
Gauduchon test limits and volume-variation sections, and the
P15 warped-cone calculation;
\url{https://github.com/MingchenXia/non-kahler-pluripotential-theory}.
\end{thebibliography}
\end{document}
